% Irodov Arihant PDF pp. 33--35, printed pp. 24--26.
\exproblem{141}{24}
如图1.35，质量 $m=1.0\,\mathrm{kg}$ 的小块置于一块
可水平缓慢移动的木板上，借一根初长
$\ell_0=0.40\,\mathrm m$ 的轻弹性绳连接到固定点 $O$。
绳初始竖直。小块与木板摩擦系数 $k=0.20$。
缓慢向右移动木板，直至弹性绳与竖直夹角
$\theta=30^\circ$、小块即将相对木板滑动。
求截至此时摩擦力对小块在地面系中所做的功。
\exsolution
设弹性绳劲度系数为 $\kappa$，
小块水平移动 $x$，绳长
$L=\sqrt{\ell_0^2+x^2}$，
伸长 $\Delta L=L-\ell_0$。
在每个准静态位置，木板摩擦力大小等于
弹性绳水平分力，
$f=\kappa(L-\ell_0)x/L$。
因 $\mathrm dL=(x/L)\mathrm dx$，
摩擦力累计做功
\[
W_f=\int_0^{x_*}f\,\mathrm dx
=\int_{\ell_0}^{L_*}\kappa(L-\ell_0)\,\mathrm dL
=\frac12\kappa(\Delta L_*)^2.
\]
到将滑动的临界位置，支持力
$N=mg-\kappa\Delta L_*\cos\theta$，而
$f=kN$，所以
\[
\kappa\Delta L_*(\sin\theta+k\cos\theta)=kmg.
\]
又 $\Delta L_*=\ell_0(\sec\theta-1)$，
故
\[
\boxed{W_f=
\frac{kmg\ell_0(1-\cos\theta)}
 {2\cos\theta(\sin\theta+k\cos\theta)}
\simeq0.090\,\mathrm J}.
\]
\excheck 此功为正，因为地面系中木板摩擦力推动小块向右；
绳储存的弹性能恰好是该做功。
板—块间接近静止的相对位移并不等于地面系的位移。

\exproblem{142}{24}
光滑轻水平杆 $AB$ 能绕端点 $A$ 的竖直轴旋转。
一小滑套质量为 $m$，以一根自然长 $\ell_0$、
劲度系数 $\kappa$ 的轻弹簧与点 $A$ 相连。
问为使系统缓慢加速到角速度 $\omega$，
外界至少应做多少功。
\exsolution
设最终稳定径向位置为 $r$。
在旋转参考系中径向力平衡：
\[
\kappa(r-\ell_0)=m\omega^2 r.
\]
令 $\eta=m\omega^2/\kappa$。
当 $0\le\eta<1$ 时
$r=\ell_0/(1-\eta)$，
弹簧伸长
$\Delta r=\ell_0\eta/(1-\eta)$。
缓慢加速可使径向振动动能趋零，
外力至少须提供最终旋转动能和弹簧弹性势能：
\[
W_{\min}
=\frac12m\omega^2r^2
+\frac12\kappa(\Delta r)^2.
\]
代入化简
\[
\boxed{W_{\min}=
\frac{\kappa\ell_0^2}{2}
\frac{\eta(1+\eta)}{(1-\eta)^2}},
\qquad
\eta=\frac{m\omega^2}{\kappa}<1.
\]
\excheck $\eta\to0$ 时
$W_{\min}\sim \tfrac12m\omega^2\ell_0^2$；
$\eta\to1^-$ 时位移与所需功发散，
意味着理想弹簧模型没有有限稳定平衡。

\exproblem{143}{24}
轻绳跨过固定的无摩擦轻滑轮，两端挂质量
$m_1,m_2$ 的重物。求两物体构成的系统质心加速度。
\exsolution
取竖直向下为正，两端加速度大小相等、方向相反。
阿特伍德机中 $m_1$ 的加速度为
\[
a_1=g\frac{m_1-m_2}{m_1+m_2},
\quad a_2=-a_1.
\]
系统总质量 $M=m_1+m_2$，
质心加速度按定义为
\[
a_C=\frac{m_1a_1+m_2a_2}{M}
=\frac{m_1-m_2}{m_1+m_2}a_1.
\]
故
\[
\boxed{a_C=g\frac{(m_1-m_2)^2}{(m_1+m_2)^2}},
\]
沿竖直向下方向（相等质量时为零）。
\excheck 虽然一个重物向上运动，另一重物向下运动，
质心仍可能向下加速；不可把两物体的加速度简单平均。

\exproblem{144}{25}
两个相互作用的质点构成孤立系统，质心始终静止。
图1.36给出了某时刻两质点的位置，
以及质量为 $m_1$ 的质点的轨迹。
另有 $m_2=m_1/2$。求画出另一个质点的轨迹。
\exsolution
由于质心静止，它的位置矢量 $\bm R_C$ 不随时间变化：
\[
m_1(\bm r_1-\bm R_C)+
m_2(\bm r_2-\bm R_C)=\bm0.
\]
利用 $m_2=m_1/2$，得到逐时刻的位置映射
\[
\boxed{\bm r_2(t)-\bm R_C
=-2[\bm r_1(t)-\bm R_C]}.
\]
因此，以固定质心为中心将图中质点1的整条轨迹
进行\emph{中心反演}，再将相对质心的距离扩大 $2$ 倍，
即得到质点2的轨迹。任意相应时刻两质点与质心共线、
位于质心的两侧。
\excheck 本版题面英文末句重复称“画 $m_1$ 的轨迹”，
而图中已给出 $m_1$ 的轨迹；
本稿按显然的物理意图给出 $m_2$ 的轨迹，
并将这一题面指代问题列为待校勘事项。

\exproblem{145}{25}
一闭合链条质量为 $m=0.36\,\mathrm{kg}$，
由一根细绳系在竖直旋转轴上，
以恒角速度 $\omega=35\,\mathrm{rad/s}$ 绕轴转动。
细绳与竖直方向夹角为 $\theta=45^\circ$（图1.37）。
求链条质心到旋转轴的距离 $r$ 及绳张力 $T$。
\exsolution
把整个闭合链条视为整体，内部约束力合力为零。
质心受力仅有向下重力 $mg$ 和沿绳向上、向轴的张力。
竖直平衡与水平向心运动分别满足
\[
T\cos\theta=mg,\qquad
T\sin\theta=m\omega^2r.
\]
两式相除消去 $T$：
\[
\boxed{r=\frac{g\tan\theta}{\omega^2}
\simeq0.00800\,\mathrm m=0.80\,\mathrm{cm}},
\]
再由竖直平衡得
\[
\boxed{T=\frac{mg}{\cos\theta}
\simeq4.99\,\mathrm N}.
\]
\excheck 只用整个链条的质心动力学便可求此二量，
不必知道链条实际受旋转影响形成的空间形状。

\exproblem{146}{25}
一均匀圆锥 $A$ 的质量
$m=3.2\,\mathrm{kg}$、半顶角 $\alpha=10^\circ$，
沿另一圆锥面 $B$ 匀速纯滚动，锥顶 $O$ 固定（图1.38）。
$A$ 的质心与 $O$ 等高，距 $O$ 的距离
$\ell=0.17\,\mathrm m$，其轴以角速度 $\omega$ 转动。
求：（a）$\omega=1.0\,\mathrm{rad/s}$ 时静摩擦力；
（b）当摩擦系数 $k=0.25$ 时保持纯滚动的角速度范围。
\exsolution
质心绕竖直轴作水平圆周运动，加速度大小为
$\omega^2\ell$，方向指向 $O$。
把接触处作用力分解成法向支持力 $N$ 与沿接触母线的
静摩擦力 $f$。在过 $O$ 和质心的竖直剖面内，
将重力与质心的水平向心加速度分别投影到接触母线
及其法线，得
\[
\boxed{f=m(g\sin\alpha+\omega^2\ell\cos\alpha)},\qquad
\boxed{N=m(g\cos\alpha-\omega^2\ell\sin\alpha)}.
\]
当 $\omega=1.0\,\mathrm{rad/s}$，
\[
\boxed{f\simeq
3.2[9.8\sin10^\circ+0.17\cos10^\circ]
\simeq5.98\,\mathrm N}.
\]
纯滚动要求 $f\le kN$，所以
\[
g\sin\alpha+\omega^2\ell\cos\alpha
\le k(g\cos\alpha-\omega^2\ell\sin\alpha).
\]
若 $k>\tan\alpha$，其角速度上限为
\[
\boxed{0\le\omega\le
\sqrt{\frac{g(k-\tan\alpha)}
{\ell(1+k\tan\alpha)}}\simeq2.02\,\mathrm{rad/s}}.
\]
\excheck 题设 $k=0.25>\tan10^\circ$，
故静止或缓慢滚动时有足够摩擦；
该上限同时保证 $N>0$。
书末给出的“$\omega\ll2$”应理解为不超过静摩擦临界值，
而非要求速度必须远小于某数值的渐近条件。

\exproblem{147}{25}
在惯性系 $K$ 中，两质点质量分别为 $m_1,m_2$，
沿同一 $x$ 轴以速率的有符号值 $v_1,v_2$ 运动。
求：（a）使两质点总动能最小的匀速平移参考系 $K'$
相对 $K$ 的速度 $V$；
（b）该最小总动能。
\exsolution
在以 $V$ 平移的惯性系中两物体的速度为
$v_1-V$ 和 $v_2-V$，总动能
\[
K(V)=\frac12m_1(v_1-V)^2+
\frac12m_2(v_2-V)^2.
\]
对 $V$ 微分求极小值：
\[
\frac{\mathrm dK}{\mathrm dV}
=-m_1(v_1-V)-m_2(v_2-V)=0.
\]
因此
\[
\boxed{V=\frac{m_1v_1+m_2v_2}{m_1+m_2}},
\]
即质心速度。代回得
\[
\boxed{K_{\min}=
\frac12\mu(v_1-v_2)^2},\quad
\mu=\frac{m_1m_2}{m_1+m_2}.
\]
\excheck 二阶导数
$\mathrm d^2K/\mathrm dV^2=m_1+m_2>0$，
故这确实是全局最小值。

\exproblem{148}{25}
某粒子系统总质量为 $M$，
其质心静止系相对惯性系 $K$ 以速度 $\bm V$ 匀速平移。
系统在质心系中的总机械能为 $\widetilde E$，
求在 $K$ 系中的总机械能 $E$。
\exsolution
质点 $i$ 的速度写为
$\bm v_i=\bm V+\widetilde{\bm v}_i$，
且质心系满足
$\sum_i m_i\widetilde{\bm v}_i=\bm0$。
因此
\[
\sum_i\frac12m_i v_i^2
=\frac12M V^2+
\sum_i\frac12m_i\widetilde v_i^2,
\]
因为动能展开后的交叉项为零。
系统内部势能与整体平移速度无关，
故总机械能关系为
\[
\boxed{E=\widetilde E+\frac12M V^2}.
\]
\excheck 这就是柯尼希动能定理；
若“总能量”包含狭义相对论静能，则不能直接将此
非相对论经典式作为严格关系。

\exproblem{149}{25--26}
两个小圆片质量分别为 $m_1,m_2$，
通过轻弹簧相连并置于光滑水平面上，
随后分别给予互相垂直的初速度矢量
$\bm v_1,\bm v_2$。
求系统在质心系中的总能量（取弹簧初始未形变）。
\exsolution
在初始时刻，系统的质心系动能为
\[
K_{\rm rel}
=\frac12\mu|\bm v_1-\bm v_2|^2,\qquad
\mu=\frac{m_1m_2}{m_1+m_2}.
\]
由于 $\bm v_1\cdot\bm v_2=0$，
\[
|\bm v_1-\bm v_2|^2=v_1^2+v_2^2.
\]
若弹簧初始处于自然长度，
则初始弹性势能为零，总能量守恒给出
\[
\boxed{\widetilde E=
\frac12\frac{m_1m_2}{m_1+m_2}(v_1^2+v_2^2)}.
\]
\excheck 若初始弹簧已发生形变，题目需要额外给出初始弹性能，
此结果应再加该项；不能由速度数据单独决定全部机械能。

\exproblem{150}{26}
两个小球质量为 $m_1,m_2$，通过轻弹簧连接。
$t=0$ 时其初速度分别为矢量
$\bm v_1,\bm v_2$，随后在地球近地面均匀重力场中运动，
忽略空气阻力。求系统总动量 $\bm p(t)$
和质心相对初始位置的位矢 $\Delta\bm r_C(t)$。
\exsolution
系统内部弹簧力总和为零，唯一外力合力是
$(m_1+m_2)\bm g$，
因此由总动量定理
\[
\frac{\mathrm d\bm p}{\mathrm dt}=(m_1+m_2)\bm g,
\]
积分得到
\[
\boxed{\bm p(t)=m_1\bm v_1+m_2\bm v_2
 +(m_1+m_2)\bm g\,t}.
\]
质心速度为
$\bm v_C(t)=\bm p(t)/(m_1+m_2)$，
所以
\[
\boxed{\Delta\bm r_C(t)=
\frac{m_1\bm v_1+m_2\bm v_2}{m_1+m_2}\,t
+\frac12\bm g\,t^2}.
\]
\excheck 无论弹簧是否振动、两球间距如何变化，
只要外力只含均匀重力，质心仍作标准匀加速抛体运动。
